\section{Introduction}
\label{sec:introduction}

\denis{} (\emph{Doppler Estimation and Numerical Inference for
Spectroscopy}) is an interactive desktop application for planning,
analysing, and reporting collinear laser spectroscopy (CLS)
measurements.

In a CLS experiment, an ion or atom beam is overlapped collinearly (or
anti-collinearly) with a frequency-narrow laser, and the hyperfine
structure (HFS) of an electronic transition is mapped by scanning the
beam acceleration voltage to Doppler-tune the laser into resonance with
each isotope. Quantitative analysis of the resulting spectra yields
isotope shifts, magnetic dipole and electric quadrupole moments, and
spin assignments --- the primary observables of nuclear structure
measurements at facilities such as IGISOL, ISOLDE, and FRIB. For
background see, e.g., the Handbook of Nuclear Physics chapter on
nuclear spins and moments~\cite{cls_handbook}.

\denis{} brings the full experiment cycle into a single tool, organised
as four top-level tabs:

\begin{description}
    \item[\app{Estimate}.]
        Predicts beam time. Given an element, transition, isotope chain
        (with masses, spins, magnetic moments and quadrupole moments),
        beam parameters, and yields, it computes per-peak signal rates,
        the time required to reach a target statistical uncertainty,
        and a total run-time budget.
    \item[\app{Pre-Analysis}.]
        Loads experiment data files (IGISOL-style ASDF via
        \code{clstools}~\cite{clstools}) and lets the user inspect the
        time-of-flight gate, set the binning axis (acceleration voltage
        or Doppler-shifted rest-frame frequency), and overlay a trial
        HFS model to seed initial fit parameters.
    \item[\app{Analysis}.]
        Performs the actual fitting. Each analysis project chains a
        \app{Source}, one or more \app{Model} blocks, a \app{Fitter},
        and one or more \app{Output} blocks. Multi-source and
        multi-model fits share parameters across runs. An
        \app{Auto-Fitter} explores the parameter space with a
        multi-start parallel \code{leastsq} sweep when the human seed
        is too far from the global minimum. Fitting itself is powered
        by \textsc{satlas2}~\cite{satlas2}.
    \item[\app{Results}.]
        Collects per-project fit summaries, exposes the matplotlib
        figures for editing (line widths, fonts, colours, axis limits)
        with persistent style files, and exports plots in PNG, PDF, or
        SVG. A permanent \app{Isotope Shifts} sub-tab combines fit
        centroids across projects to compute $\delta\nu$ relative to a
        chosen reference isotope.
\end{description}

A \menu{Tools} menu provides supporting utilities (Schmidt moment
calculator, shell configuration plotter, unit converter, SHG crystal
angle calculator, SFG / DFG calculator, quick plot, and a NIST
Atomic Spectra Database browser with a laser excitation-scheme
finder) for tasks adjacent to the main workflow; see
Section~\ref{sec:tools}.

\paragraph{Supported data format.}
The \app{Pre-Analysis} and \app{Analysis} tabs in the current release
read \emph{only} IGISOL-style ASDF files, via the
\code{clstools}~\cite{clstools} reader. Other facility formats are not
supported out of the box; adapting \denis{} to a different layout
requires editing the source-block loader. The \app{Estimate} tab has
no file-format dependency and can be used independently of any data
on disk.

\paragraph{Audience.}
This manual assumes the reader has hands-on experience running CLS
measurements and is comfortable with concepts such as Doppler tuning,
hyperfine structure, the $A$ and $B$ coupling constants, and Schmidt
moments. For background on the underlying physics see
References~\cite{satlas2,cls_handbook} and
Appendix~\ref{sec:physics}.

\paragraph{Scope.}
This is a user manual, not an API reference: it documents the GUI
workflow and supporting concepts, not the internal Python modules. The
backend code (in \code{cls\_estimations/} and \code{gui/}) is intended
to be read directly by developers extending \denis{}.
